\emph{(i)} For a standard bivariate Gaussian pair $(X, Y)$ with correlation $\rho$, Mehler's formula (Fact \claimref{synthesis.imported-mehler-formula-hermite-rank-and-power-law-spectra}) gives $E[\mathrm{He}_{m}(X)\mathrm{He}_{n}(Y)] = \delta_{mn}\, n!\, \rho^{n}$. A proof: the generating function $e^{tx - t^{2}/2} = \sum_{n} \mathrm{He}_{n}(x)t^{n}/n!$ gives $E[e^{tX - t^{2}/2}e^{uY - u^{2}/2}] = e^{tu\rho}$ from the Gaussian moment generating function, and comparing coefficients of $t^{m}u^{n}$ gives the identity. The check verifies it exactly for $m, n \le 7$ as polynomials in $\rho$, computing every joint moment from $Y = \rho X + (1 - \rho^{2})^{1/2}Z$ with $Z$ independent (only even powers of the square root survive). For the field, $(x(\mathbf{x}), x(\mathbf{y}))$ is such a pair with $\rho = \rho(r)$; for a quasi-free state the same holds for Wick powers at spacelike separation, where the two-point function is real and Wick's theorem applies. Since $f(\sigma)$ has finite variance, $\sum c_{n}^{2}n! < \infty$, the expansion converges in mean square and $\mathrm{Cov}(\zeta(\mathbf{x}), \zeta(\mathbf{y})) = \sum_{n \ge 1}c_{n}^{2}n!\rho(r)^{n}$, the $n = 0$ term being the mean. The check also computes the covariance directly from moments for the dominant curvaton's series $\tfrac23 x - \tfrac13 x^{2} + \tfrac29 x^{3}$ (Fact \claimref{synthesis.imported-curvaton-and-modulated-reheating-coefficients}), for a modulated decay rate $\Gamma \propto \sigma^{3}$ read through $-\tfrac16 \ln\Gamma$ to fourth order, and for a pure square, and finds the Hermite form with non-negative coefficients and lowest power $n_{0}$ in each.

\emph{(ii)} For $0 < \rho \le 1$ the series $g(\rho) = \sum_{n \ge n_{0}}c_{n}^{2}n!\rho^{n}$ converges with positive terms, and $\rho g'(\rho) = \sum_{n}n\,c_{n}^{2}n!\rho^{n}$ (termwise differentiation inside the radius of convergence; at $\rho = 1$ as a limit, possibly infinite). Hence $\bar n = \rho g'/g$ is the mean of $n$ under the weights $c_{n}^{2}n!\rho^{n}$, at least $n_{0}$; if $c_{1} \neq 0$, $\bar n - 1 = \sum_{n \ge 2}(n - 1)c_{n}^{2}n!\rho^{n}/g = O(\rho)$ as $\rho \to 0$. The chain rule gives $d\ln g(\rho(r))/d\ln r = \bar n\, d\ln\rho/d\ln r$. The check verifies $\rho g'/g \ge n_{0}$ exactly at sample $\rho \in (0, 1]$ for the three conversions of (i), and the bound $\bar n - 1 \le \rho\,\sum_{n \ge 2}(n - 1)c_{n}^{2}n!/c_{1}^{2}$ for $0 < \rho \le 1$ for the curvaton series, where the constant is $11/24$.

\emph{(iii)} With the power-law hypothesis, since $|\rho| \le 1$ and $\rho(r) \to 0$, for $r$ large $\mathrm{Cov} = c_{n_{0}}^{2}n_{0}!\,A^{n_{0}}r^{-n_{0}m}(1 + o(1))$, the higher terms being $O(\rho^{n_{0}+1})$ uniformly by the convergent sum. The Fourier transform of a correlation that behaves as $B r^{-\beta}$ at large $r$ (and is locally integrable) behaves, for $0 < \beta < 3$, as $P(k) \sim 4\pi B\,\Gamma(2 - \beta)\sin(\tfrac{\pi}{2}(2 - \beta))\,k^{\beta - 3}$ as $k \to 0$ (Fact \claimref{synthesis.imported-mehler-formula-hermite-rank-and-power-law-spectra}; the weight is positive on $0 < \beta < 3$, as the check's sign table by intervals confirms), so the dimensionless spectrum $k^{3}P(k)/2\pi^{2}$ behaves as $k^{\beta}$. With $\beta = n_{0}m < 3$ the index is $n_{0}m$. If $n_{0}m > 3$ the covariance is integrable, $P(0) = \int \mathrm{Cov}\,d^{3}r$ is finite and non-negative (a power spectrum is non-negative), and the dimensionless spectrum is $P(0)k^{3}/2\pi^{2}$ plus higher powers: index $3$, or more if $P(0) = 0$. In both cases the index is at least $m$, with equality when $n_{0} = 1$. The check verifies $\min(n_{0}m, 3) \ge m$ on a grid of $m$ in $(0, 3)$ and $n_{0} \le 6$.

\emph{(iv)} Theorem \claimref{synthesis.frozen-point-geometry-reads-descent-as-vacuum} gives, under hypothesis (F), the stated indices of every stabiliser-invariant Hadamard state on the descent slice: exactly $2$ for the conformal weight (correlator $1/(4\pi^{2}a^{2}r^{2})$), $2u^{2}/(1 + u^{2})$ for the massless Bunch-Davies state ($\Delta^{2} = (H/2\pi)^{2}(1 + u^{2})$, $u = k|x_{d}|$), and between $2 - 2/(u^{2} - 2)$ and $3$ for every unitary-series field at leading adiabatic order, with $u > 2.7$ at the largest observed scales and $u \in (960, 968)$ at the pivot. At $u = 27/10$ these are at least $1458/829 > 1.75$ (massless) and $2 - 200/529 > 1.62$ (unitary series); at $u = 960$ both exceed $1.99$. Those states are quasi-free, so (i) applies to the field smeared over the descent's Hubble length. At observed separations $r$, which exceed that length by more than $2 \times 10^{5}$, the smeared field's normalised correlation is positive and small (argued: the vacuum part is of order the squared ratio of the smearing length to $r$, below $10^{-10}$, and the de Sitter part is suppressed by the squared ratio of the frozen point's Hubble rate to the descent's), so by (ii) a conversion with $c_{1} \neq 0$ keeps the field's local index up to a relative correction of that order, and one with $c_{1} = 0$ at least doubles its local rate of fall. Against $n_{s} - 1 = -0.0351 \pm 0.0042$ (Fact \claimref{synthesis.imported-planck-2018-lcdm-parameters}) the gaps are more than $(1.62 + 0.0351)/0.0042 > 394$ standard deviations at the largest observed scales and more than $(1.99 + 0.0351)/0.0042 > 482$ at the pivot. The modulation of the descent's deposit by a field, the non-quasi-free escape recorded in Theorem \claimref{synthesis.quasi-free-deposit-is-vacuum-plus-shot-noise}, is a local conversion of this kind when the modulating field is Gaussian, so it inherits the bound. A red curvature perturbation from such a conversion needs a modulating field that is already red. The general lemma (i) to (iii) is the classical Hermite-rank result (Taqqu 1975; Dobrushin and Major 1979; Fact \claimref{synthesis.imported-mehler-formula-hermite-rank-and-power-law-spectra}); the series enters only through (iv).
